% Level 1, Week 6 -- workbook chapter
\chapterweek{6}{Review Studio and Checkpoint}{Show what you know; get ready for text code.}

\section{What you will learn}
\begin{itemize}[leftmargin=1.2em,itemsep=1pt]
  \item How Scratch blocks turn into \textbf{Python} words (next week's language!).
  \item How to be a bug detective in projects someone else made.
  \item What makes a project \textbf{fun} and \textbf{fair} for everyone.
\end{itemize}

\section{The big idea: our six big ideas}
Every programming language in the world uses these six ideas. You know them all!

\begin{center}
\renewcommand{\arraystretch}{1.5}
\begin{tabular}{@{}m{2.6cm}m{6.4cm}m{6.4cm}@{}}
\textbf{Big idea} & \textbf{It means\ldots} & \textbf{In Scratch} \\ \hline
\textbf{Sequence} (week 1) & Steps happen in order, one after another. & \begin{scratch}[scale=0.6]\blockmove{move \ovalnum{100} steps}\blocklook{say \ovalnum{Hi!}}\end{scratch} \\
\textbf{Loops} (week 2) & Repeat steps so you don't copy them. & \begin{scratch}[scale=0.6]\blockrepeat{repeat \ovalnum{4}}{\blockmove{move \ovalnum{100} steps}}\end{scratch} \\
\textbf{Events} (week 3) & A doorbell starts a script. & \begin{scratch}[scale=0.6]\blockinit{when \selectmenu{space} key pressed}\end{scratch} \\
\textbf{Conditionals} (week 3) & Ask a yes/no question, then choose. & \begin{scratch}[scale=0.6]\blockif{if \boolsensing{touching \ovalsensing*{edge} ?} then}{\blockspace}\end{scratch} \\
\textbf{Variables} (week 4) & A labeled box that remembers. & \begin{scratch}[scale=0.6]\blockvariable{change \selectmenu{score} by \ovalnum{1}}\end{scratch} \\
\textbf{My Blocks} (week 5) & Name a part and use it again. & \begin{scratch}[scale=0.6]\blockmoreblocks{polygon \ovalnum{6} \ovalnum{60}}\end{scratch} \\
\end{tabular}
\end{center}

\begin{tip}[The four detective steps]
\textbf{1.} What did I \textbf{expect}? \quad \textbf{2.} What \textbf{happened}? \quad \textbf{3.} Read the blocks one at a time with your finger. \quad \textbf{4.} Change \textbf{one} thing, then test again.
\end{tip}

\section{Bug-hunt relay}
Each project card shows a project that someone built with \textbf{two bugs}. Build it \textbf{exactly} as shown (bugs and all!), click the flag, then be detectives. Fill in the relay baton for each one.

\begin{exercise}[Relay A: Wonky Hexagon]
\textbf{It should:} draw one closed hexagon (6 equal sides of 80 steps) in the middle of the stage, then the cat says \enquote{Done!}.

\textbf{Build it:} new project, add the \menu{Pen} extension, set the cat's \menu{Size} to 40, then build this script on the cat:
\begin{center}
\begin{scratch}[scale=0.78]
\blockinit{when \greenflag clicked}
\blockpen{erase all}
\blockmove{go to x: \ovalnum{0} y: \ovalnum{0}}
\blockmove{point in direction \ovalnum{90}}
\blockpen{pen down}
\blockrepeat{repeat \ovalnum{6}}{
  \blockmove{move \ovalnum{80} steps}
  \blockmove{turn \turnright{} \ovalnum{50} degrees}
  \blockpen{pen up}
}
\blocklook{say \ovalnum{Done!} for \ovalnum{2} seconds}
\end{scratch}
\end{center}
\textbf{Relay baton.} What happened? \answerlines{1}
Bug 1 and fix: \answerlines{1}
Bug 2 and fix: \answerlines{1}
\end{exercise}

\begin{exercise}[Relay B: Apple Catch Can't Count]
\textbf{It should:} start every game with a score of 0; drop apples from surprise places at the top; add \textbf{exactly 1} for every apple caught in the bowl; send missed apples back to the top.

\textbf{Build it:} new project, delete the cat, add the \textbf{Bowl} and \textbf{Apple} sprites, make a variable \texttt{score} (for all sprites), then:
\begin{center}
\begin{tabular}{@{}c@{\hspace{0.8cm}}c@{}}
\begin{scratch}[scale=0.72]
\blockinit{when \greenflag clicked}
\blockmove{go to x: \ovalnum{0} y: \ovalnum{-140}}
\blockinfloop{forever}{
  \blockmove{set x to \ovalsensing{mouse x}}
}
\end{scratch}
&
\begin{scratch}[scale=0.72]
\blockinit{when \greenflag clicked}
\blockmove{go to x: \ovaloperator{pick random \ovalnum{-200} to \ovalnum{200}} y: \ovalnum{170}}
\blockinfloop{forever}{
  \blockmove{change y by \ovalnum{-5}}
  \blockif{if \boolsensing{touching \ovalsensing*{Bowl} ?} then}{
    \blockvariable{change \selectmenu{score} by \ovalnum{1}}
  }
  \blockif{if \booloperator{\ovalmove{y position} < \ovalnum{-170}} then}{
    \blockmove{go to x: \ovaloperator{pick random \ovalnum{-200} to \ovalnum{200}} y: \ovalnum{170}}
  }
}
\end{scratch}
\\[4pt]
{\small\textbf{Bowl}} & {\small\textbf{Apple}}
\end{tabular}
\end{center}
Play two games. \textbf{Relay baton.} What happened? \answerlines{1}
Bug 1 and fix: \answerlines{1}
Bug 2 and fix: \answerlines{1}
\end{exercise}

\begin{exercise}[Relay C: The Silent Party]
\textbf{It should:} after the green flag, the Ghost hides. When you press \keys{space}, the cat says \enquote{Party time!}, then the Ghost appears, does a slow spinning dance you can really watch (one full turn, about 2 seconds) and says \enquote{Woo!}.

\textbf{Build it:} new project, keep the cat, add the \textbf{Ghost}. Make two messages called \texttt{party} and \texttt{dance} (broadcast dropdown $\rightarrow$ \menu{New message}). Then:
\begin{center}
\begin{tabular}{@{}c@{\hspace{0.8cm}}c@{}}
\begin{scratch}[scale=0.72]
\blockinit{when \selectmenu{space} key pressed}
\blocklook{say \ovalnum{Party time!} for \ovalnum{2} seconds}
\blockevent{broadcast \selectmenu{party}}
\end{scratch}
&
\begin{scratch}[scale=0.72]
\blockinit{when \greenflag clicked}
\blocklook{hide}
\end{scratch}
\\[6pt]
&
\begin{scratch}[scale=0.72]
\blockinit{when I receive \selectmenu{dance}}
\blocklook{show}
\blockrepeat{repeat \ovalnum{10}}{
  \blockmove{turn \turnright{} \ovalnum{36} degrees}
  \blocklook{next costume}
}
\blocklook{say \ovalnum{Woo!} for \ovalnum{2} seconds}
\end{scratch}
\\[4pt]
{\small\textbf{Cat}} & {\small\textbf{Ghost}}
\end{tabular}
\end{center}
\textbf{Relay baton.} What happened? \answerlines{1}
Bug 1 and fix: \answerlines{1}
Bug 2 and fix: \answerlines{1}
\end{exercise}

\section{Blocks to words}
Next week you will write code in \textbf{Python}. Python uses words instead of blocks. Match each block card (A--L) with its Python card (1--12). Write the number next to each letter.

\begin{tip}[Clues for the card sort]
\begin{itemize}[leftmargin=1.2em,itemsep=1pt]
  \item In Python we draw with a \textbf{turtle}: a little arrow that carries a pen, like the Scratch cat with the Pen extension. In our code the turtle is called \texttt{t}. So a line that starts with \texttt{t.} is an order for the turtle.
  \item Look for a word you know. Many Python words sound like the words on the blocks.
  \item Guessing is fine! We check the answers together.
\end{itemize}
\end{tip}

\begin{center}
\renewcommand{\arraystretch}{1.6}
\begin{tabular}{|*{12}{c|}}
\hline
\textbf{A} & \textbf{B} & \textbf{C} & \textbf{D} & \textbf{E} & \textbf{F} & \textbf{G} & \textbf{H} & \textbf{I} & \textbf{J} & \textbf{K} & \textbf{L} \\ \hline
\rule{0.75cm}{0pt} & \rule{0.75cm}{0pt} & \rule{0.75cm}{0pt} & \rule{0.75cm}{0pt} & \rule{0.75cm}{0pt} & \rule{0.75cm}{0pt} & \rule{0.75cm}{0pt} & \rule{0.75cm}{0pt} & \rule{0.75cm}{0pt} & \rule{0.75cm}{0pt} & \rule{0.75cm}{0pt} & \rule{0.75cm}{0pt} \\ \hline
\end{tabular}
\end{center}
What do you notice about Python? (Look at the parentheses, quotes and colons.) \answerlines{1}

\begin{vocab}
  \vocabword{debugging}{Finding and fixing bugs, step by step, like a detective.}
  \vocabword{Python}{A programming language where you type words instead of snapping blocks.}
  \vocabword{text code}{Code made of typed words and symbols.}
  \vocabword{checkpoint}{A short check that shows what you know and what to practice next.}
  \vocabword{fair}{Same rules for everyone, no tricks, and everyone can play.}
\end{vocab}

\begin{reflection}[Fun and fair for everyone]
One thing that makes a game \textbf{fun}: \answerlines{1}
One thing that makes a game \textbf{fair}: \answerlines{1}
The project I am proudest of from weeks 1--6, and why: \answerlines{2}
\end{reflection}

\clearpage
\subsection{Printable: block cards (cut along the lines)}
\begin{center}
\setlength{\fboxsep}{0pt}
\fbox{\begin{minipage}[c][2.6cm][c]{0.31\linewidth}\centering{\large\bfseries A}\\[4pt]
\begin{scratch}[scale=0.5]\blockrepeat{repeat \ovalnum{4}}{\blockspace}\end{scratch}\end{minipage}}%
\fbox{\begin{minipage}[c][2.6cm][c]{0.31\linewidth}\centering{\large\bfseries B}\\[4pt]
\begin{scratch}[scale=0.8]\blocklook{say \ovalnum{Hello!}}\end{scratch}\end{minipage}}%
\fbox{\begin{minipage}[c][2.6cm][c]{0.31\linewidth}\centering{\large\bfseries C}\\[4pt]
\begin{scratch}[scale=0.8]\blockpen{pen up}\end{scratch}\end{minipage}}

\fbox{\begin{minipage}[c][2.6cm][c]{0.31\linewidth}\centering{\large\bfseries D}\\[4pt]
\begin{scratch}[scale=0.8]\blockvariable{change \selectmenu{score} by \ovalnum{1}}\end{scratch}\end{minipage}}%
\fbox{\begin{minipage}[c][2.6cm][c]{0.31\linewidth}\centering{\large\bfseries E}\\[4pt]
\begin{scratch}[scale=0.8]\blockmove{move \ovalnum{100} steps}\end{scratch}\end{minipage}}%
\fbox{\begin{minipage}[c][2.6cm][c]{0.31\linewidth}\centering{\large\bfseries F}\\[4pt]
\begin{scratch}[scale=0.6]\blockvariable{set \selectmenu{roll} to \ovaloperator{pick random \ovalnum{1} to \ovalnum{6}}}\end{scratch}\end{minipage}}

\fbox{\begin{minipage}[c][2.6cm][c]{0.31\linewidth}\centering{\large\bfseries G}\\[4pt]
\begin{scratch}[scale=0.8]\blockmove{turn \turnright{} \ovalnum{90} degrees}\end{scratch}\end{minipage}}%
\fbox{\begin{minipage}[c][2.6cm][c]{0.31\linewidth}\centering{\large\bfseries H}\\[4pt]
\begin{scratch}[scale=0.8]\blockvariable{set \selectmenu{score} to \ovalnum{0}}\end{scratch}\end{minipage}}%
\fbox{\begin{minipage}[c][2.6cm][c]{0.31\linewidth}\centering{\large\bfseries I}\\[4pt]
\begin{scratch}[scale=0.8]\blockmove{go to x: \ovalnum{0} y: \ovalnum{0}}\end{scratch}\end{minipage}}

\fbox{\begin{minipage}[c][2.6cm][c]{0.31\linewidth}\centering{\large\bfseries J}\\[4pt]
\begin{scratch}[scale=0.5]\blockif{if \booloperator{\ovalvariable{score} > \ovalnum{10}} then}{\blockspace}\end{scratch}\end{minipage}}%
\fbox{\begin{minipage}[c][2.6cm][c]{0.31\linewidth}\centering{\large\bfseries K}\\[4pt]
\begin{scratch}[scale=0.8]\blockpen{pen down}\end{scratch}\end{minipage}}%
\fbox{\begin{minipage}[c][2.6cm][c]{0.31\linewidth}\centering{\large\bfseries L}\\[4pt]
\begin{scratch}[scale=0.75]\blocksensing{ask \ovalnum{Nickname?} and wait}\end{scratch}\end{minipage}}
\end{center}

\clearpage
\subsection{Printable: Python cards (cut along the lines)}
\begin{center}
\setlength{\fboxsep}{0pt}
\fbox{\begin{minipage}[c][2.6cm][c]{0.31\linewidth}\centering{\large\bfseries 1}\\[6pt]\texttt{\large t.pendown()}\end{minipage}}%
\fbox{\begin{minipage}[c][2.6cm][c]{0.31\linewidth}\centering{\large\bfseries 2}\\[6pt]\texttt{\large score = 0}\end{minipage}}%
\fbox{\begin{minipage}[c][2.6cm][c]{0.31\linewidth}\centering{\large\bfseries 3}\\[6pt]\texttt{\large for i in range(4):}\end{minipage}}

\fbox{\begin{minipage}[c][2.6cm][c]{0.31\linewidth}\centering{\large\bfseries 4}\\[6pt]\texttt{roll =}\\\texttt{random.randint(1, 6)}\end{minipage}}%
\fbox{\begin{minipage}[c][2.6cm][c]{0.31\linewidth}\centering{\large\bfseries 5}\\[6pt]\texttt{\large t.right(90)}\end{minipage}}%
\fbox{\begin{minipage}[c][2.6cm][c]{0.31\linewidth}\centering{\large\bfseries 6}\\[6pt]\texttt{\large print("Hello!")}\end{minipage}}

\fbox{\begin{minipage}[c][2.6cm][c]{0.31\linewidth}\centering{\large\bfseries 7}\\[6pt]\texttt{\large if score > 10:}\end{minipage}}%
\fbox{\begin{minipage}[c][2.6cm][c]{0.31\linewidth}\centering{\large\bfseries 8}\\[6pt]\texttt{\large t.goto(0, 0)}\end{minipage}}%
\fbox{\begin{minipage}[c][2.6cm][c]{0.31\linewidth}\centering{\large\bfseries 9}\\[6pt]\texttt{name = input(}\\\texttt{"Nickname? ")}\end{minipage}}

\fbox{\begin{minipage}[c][2.6cm][c]{0.31\linewidth}\centering{\large\bfseries 10}\\[6pt]\texttt{\large t.penup()}\end{minipage}}%
\fbox{\begin{minipage}[c][2.6cm][c]{0.31\linewidth}\centering{\large\bfseries 11}\\[6pt]\texttt{\large score = score + 1}\end{minipage}}%
\fbox{\begin{minipage}[c][2.6cm][c]{0.31\linewidth}\centering{\large\bfseries 12}\\[6pt]\texttt{\large t.forward(100)}\end{minipage}}
\end{center}
